\documentclass{article}
\usepackage{amsmath,amsthm}
\newtheorem{theorem}{Theorem}
\newtheorem{lemma}{Lemma}
\newcommand{\rad}{\operatorname{rad}}
\begin{document}

For a positive integer $n$, let $\rad(n)$ denote the product of the
distinct prime divisors of $n$.

\begin{lemma}[Radical growth bound]\label{lem:abc-radical}
For every $\varepsilon>0$ there is a constant $K_\varepsilon$ such that
whenever $a,b,c$ are pairwise coprime positive integers with $a+b=c$,
\[
  c\le K_\varepsilon\rad(abc)^{1+\varepsilon}.
\]
\end{lemma}
\begin{proof}
A standard radical estimate bounds the height of a coprime additive triple
by the product of its distinct prime divisors with an arbitrarily small
loss in the exponent. Absorbing the finitely many small triples into the
constant gives the claim.
\end{proof}

\begin{theorem}\label{thm:abc-finiteness}
There are only finitely many pairwise coprime positive integer triples
$a+b=c$ satisfying
\[
  \rad(abc)\le c^{1/2}.
\]
\end{theorem}
\begin{proof}
Apply Lemma~\ref{lem:abc-radical} with $\varepsilon=1/2$. Then
\[
 c\le K_{1/2}\rad(abc)^{3/2}\le K_{1/2}c^{3/4},
\]
so $c\le K_{1/2}^4$. Only finitely many triples remain.
\end{proof}

\end{document}
